\clearpage\thispagestyle{empty}
\pdfbookmark[0]{Quellen zur Besprechung von Bouvard}{bouvard1808-sources}
\begin{center}\textbf{Quellen zur Besprechung von Bouvard}\end{center}
\noindent Diese KI-Übertragung ins heutige Deutsch folgt Gauß’ Besprechung auf S.~508--509 der \href{https://gdz.sub.uni-goettingen.de/id/PPN236007467}{\textit{Werke}, Bd.~VI} (Göttingen, 1874). Sie erschien ohne Namensnennung unter der Überschrift \emph{Paris} in den \textit{Göttingischen gelehrten Anzeigen}, Stück~177, 5. November 1808, S.~\href{https://api.digitale-sammlungen.de/iiif/image/v2/bsb10538674_00745/full/1800,/0/default.jpg}{1761}--\href{https://api.digitale-sammlungen.de/iiif/image/v2/bsb10538674_00748/full/1800,/0/default.jpg}{1764}. Das digitalisierte Exemplar stammt aus der Bayerischen Staatsbibliothek.
\par\smallskip\noindent Der \href{https://www.e-rara.ch/zut/content/titleinfo/21914082}{Band von 1808} wurde vom französischen Bureau des longitudes herausgegeben und bei Courcier gedruckt. Gauß bespricht Alexis Bouvards \emph{Neue Tafeln für Jupiter und Saturn, berechnet nach der Theorie von Laplace und nach der Dezimaleinteilung des rechten Winkels}. Der Band enthält auch die Tafeln der Jupitermonde von Jean-Baptiste Delambre. Das verlinkte Exemplar stammt aus der ETH-Bibliothek (Rar 15631).
\par\smallskip\noindent Der Erstdruck nennt auf \href{https://api.digitale-sammlungen.de/iiif/image/v2/bsb10538674_00746/full/1800,/0/default.jpg}{S.~1762} bei der Angabe zu S.~337 von Laplaces \textit{Mécanique céleste} den 4. Teil. Die Werkausgabe nennt den 8. Teil und hat unmittelbar vor \emph{Tafeln} ein apostrophförmiges Zeichen. Die Übertragung folgt der Werkausgabe. KI-Übertragung und KI-Satz.
\clearpage
